% ===========================================================
%  RelayGuard  --  Complete Project Development Guide
%  Supersedes: Relay_Health_System_Dev_Guide.pdf
%              RUL_Predictive_Maintenance_Solution.pdf
%              ML_Model.pdf
%  Version 3.0  |  April 2026
% ===========================================================
\documentclass[12pt,a4paper]{article}

\usepackage[T1]{fontenc}
\usepackage{lmodern}
\usepackage{geometry}
\geometry{margin=2.4cm, headheight=15pt}
\usepackage{amsmath,amssymb}
\usepackage{booktabs}
\usepackage{tabularx}
\usepackage{longtable}
\usepackage{listings}
\usepackage{xcolor}
\usepackage{graphicx}
\usepackage{hyperref}
\usepackage{enumitem}
\usepackage{fancyhdr}
\usepackage{titlesec}
\usepackage{tcolorbox}
\usepackage{array}
\usepackage{float}
\usepackage{mdframed}
\usepackage{multicol}
\usepackage{tikz}
\usepackage{microtype}
\setlength{\emergencystretch}{3em}

% ---- Colours -----------------------------------------------
\definecolor{accent}{RGB}{25,84,210}
\definecolor{accent2}{RGB}{0,137,123}
\definecolor{notebg}{RGB}{227,242,253}
\definecolor{noteln}{RGB}{25,84,210}
\definecolor{warnbg}{RGB}{255,243,224}
\definecolor{warnln}{RGB}{230,81,0}
\definecolor{fixbg}{RGB}{232,245,233}
\definecolor{fixln}{RGB}{46,125,50}
\definecolor{codebg}{RGB}{248,248,248}
\definecolor{gray1}{RGB}{97,97,97}

% ---- Code listings -----------------------------------------
\lstset{
  basicstyle      = \ttfamily\footnotesize,
  breaklines      = true,
  frame           = single,
  framesep        = 3pt,
  rulecolor       = \color{gray!35},
  backgroundcolor = \color{codebg},
  keywordstyle    = \color{accent}\bfseries,
  commentstyle    = \color{gray1}\itshape,
  stringstyle     = \color{accent2},
  numbers         = left,
  numberstyle     = \tiny\color{gray},
  numbersep       = 7pt,
  xleftmargin     = 1.5em,
  tabsize         = 4,
  extendedchars   = false,
}

% ---- Section styles ----------------------------------------
\titleformat{\section}
  {\large\bfseries\color{accent}}{\thesection.}{0.6em}{}[\color{accent}\titlerule]
\titleformat{\subsection}
  {\normalsize\bfseries\color{accent2}}{\thesubsection}{0.5em}{}
\titleformat{\subsubsection}
  {\small\bfseries}{\thesubsubsection}{0.4em}{}

% ---- Header / footer ---------------------------------------
\pagestyle{fancy}\fancyhf{}
\fancyhead[L]{\small\textbf{RelayGuard} -- Complete Project Development Guide v3.0}
\fancyhead[R]{\small April 2026}
\fancyfoot[C]{\small\thepage}
\renewcommand{\headrulewidth}{0.4pt}

% ---- Custom boxes ------------------------------------------
\newenvironment{notebox}{%
  \begin{mdframed}[backgroundcolor=notebg,linecolor=noteln,linewidth=1.5pt,
    innerleftmargin=9pt,innerrightmargin=9pt,innertopmargin=5pt,innerbottommargin=5pt]
  \textbf{\color{noteln}Note: }}{\end{mdframed}}

\newenvironment{warnbox}{%
  \begin{mdframed}[backgroundcolor=warnbg,linecolor=warnln,linewidth=1.5pt,
    innerleftmargin=9pt,innerrightmargin=9pt,innertopmargin=5pt,innerbottommargin=5pt]
  \textbf{\color{warnln}Warning: }}{\end{mdframed}}

\newenvironment{fixbox}{%
  \begin{mdframed}[backgroundcolor=fixbg,linecolor=fixln,linewidth=2pt,
    innerleftmargin=9pt,innerrightmargin=9pt,innertopmargin=6pt,innerbottommargin=6pt]
  \textbf{\color{fixln}Key Design Decision: }}{\end{mdframed}}

\newenvironment{checkgate}{%
  \begin{mdframed}[backgroundcolor=yellow!10,linecolor=orange!70,linewidth=1.5pt,
    innerleftmargin=9pt,innerrightmargin=9pt,innertopmargin=5pt,innerbottommargin=5pt]
  \textbf{\color{orange!70!black}Verification Gate (must pass before next stage): }}{\end{mdframed}}

% ---- Hyperref ----------------------------------------------
\hypersetup{colorlinks=true,linkcolor=accent,urlcolor=accent2}

% ---- Title page --------------------------------------------
\title{%
  {\LARGE\bfseries\color{accent}RelayGuard}\\[5pt]
  {\large Complete Project Development Guide}\\[7pt]
  {\normalsize\color{gray}
    Predictive Maintenance System for Industrial Relays\\[2pt]
    IoT Hardware \(\cdot\) ESP32 Firmware \(\cdot\) Raspberry Pi Gateway \(\cdot\) ML Pipeline \(\cdot\) Dashboard}
}
\author{RelayGuard Engineering Team}
\date{Version 3.0\quad|\quad April 2026}

% ============================================================
\begin{document}
% ============================================================
\maketitle
\thispagestyle{empty}

\vspace{0.5em}
\begin{tcolorbox}[colback=notebg,colframe=accent,title={\bfseries About This Document}]
This is the \textbf{single authoritative development guide} for the RelayGuard project. It supersedes all
previous documents (\textit{Relay\_Health\_System\_Dev\_Guide.pdf},
\textit{RUL\_Predictive\_Maintenance\_Solution.pdf}, \textit{ML\_Model.pdf}).

The guide covers every layer of the system -- from soldering the first sensor to deploying the trained ML
model -- in build order, with verification checkpoints, exact pin assignments, code snippets, MODBUS register
maps, MQTT topic design, and the complete updated $\Delta$\,HI machine-learning architecture.

\textbf{Primary ML change from v1/v2}: the model now predicts \textit{change in Health Index}
($\Delta$\,HI) instead of absolute future HI, eliminating covariate shift and producing accurate predictions
across the full relay lifecycle.
\end{tcolorbox}

\vspace{0.5em}
\tableofcontents
\clearpage

% ============================================================
\section{Project Vision and Architecture}
% ============================================================

\subsection{What RelayGuard Does}

RelayGuard is an IoT predictive maintenance system that monitors the health of industrial relays in real time
and predicts their Remaining Useful Life (RUL) before failure occurs. The system achieves this by:

\begin{enumerate}[noitemsep]
  \item Measuring five physical signals from each relay (current, voltage, temperature, contact resistance, switching cycles).
  \item Converting those signals into a single \textbf{Health Index (HI)} using physics-based degradation models running on the ESP32.
  \item Transmitting HI and raw readings to a Raspberry Pi gateway over MODBUS/RS485.
  \item Forwarding data to the cloud over MQTT, where an XGBoost ML model predicts future HI.
  \item Displaying fleet health on a live React dashboard and auto-generating maintenance work orders.
\end{enumerate}

\subsection{Five-Layer Architecture}

\begin{table}[H]
\centering
\renewcommand{\arraystretch}{1.35}
\begin{tabularx}{\textwidth}{llX}
\toprule
\textbf{Layer} & \textbf{Hardware / Stack} & \textbf{Responsibility} \\
\midrule
L1 -- Sensing      & ACS712, ZMPT101B, MLX90614, INA219 & Read current, voltage, temperature, contact resistance \\
L2 -- Edge Node    & ESP32-S3, FreeRTOS & Compute HI, drive LEDs/buzzer, expose MODBUS slave \\
L3 -- Gateway      & Raspberry Pi 4 & Poll MODBUS, store SQLite, run ML inference, MQTT broker \\
L4 -- Cloud        & MQTT broker, InfluxDB / SQLite & Receive data, run fleet-wide XGBoost model, serve API \\
L5 -- Dashboard    & React + Recharts (or Grafana) & Fleet overview, HI timeline, work orders, RUL charts \\
\bottomrule
\end{tabularx}
\caption{Five-layer system architecture.}
\end{table}

\subsection{End-to-End Data Flow}

\[
\underbrace{\text{Sensors}}_{\text{L1}}
\xrightarrow{\text{ADC/I2C}}
\underbrace{\text{ESP32}}_{\text{L2}}
\xrightarrow{\text{MODBUS}}
\underbrace{\text{Rasp.~Pi}}_{\text{L3}}
\xrightarrow{\text{MQTT/TLS}}
\underbrace{\text{Cloud+ML}}_{\text{L4}}
\xrightarrow{\text{HTTPS/WS}}
\underbrace{\text{Dashboard}}_{\text{L5}}
\]

\noindent OTA updates flow in the reverse direction: the cloud pushes refined degradation thresholds back to
the Pi, which relays them to each ESP32 node.

\subsection{Build Sequence at a Glance}

\begin{table}[H]
\centering
\begin{tabularx}{\textwidth}{llX}
\toprule
\textbf{Day} & \textbf{Goal} & \textbf{Go/No-Go Gate} \\
\midrule
0 & Environment setup  & Blink sketch uploads. Python + pip libraries import. Git commit pushed. \\
1 & Sensor wiring      & Each sensor reads within 5\% of reference instrument. \\
2 & ESP32 firmware     & OLED shows live HI. LEDs change as test inputs change. NVS survives reboot. \\
3 & MODBUS             & QModMaster reads all 14 registers. Values match OLED. \\
4 & Pi gateway         & SQLite grows every 5\,s. All node values present and sane. \\
5 & MQTT + API         & mosquitto\_sub shows JSON every 5\,s. /api/fleet returns JSON in browser. \\
6 & Dashboard          & HI gauge updates without page refresh. Alarm state visible within 10\,s. \\
7 & ML + integration   & ml\_inference.py runs on Pi. Anomaly flag in API. Full end-to-end test. \\
\bottomrule
\end{tabularx}
\caption{Day-by-day build sequence with mandatory verification gates.}
\end{table}

% ============================================================
\section{Stage 0 -- Environment Setup (Day 0)}
% ============================================================

\begin{notebox}
Golden rule: \textbf{build layer by layer and verify each layer independently}. A sensor bug and a MODBUS
bug look identical when they are stacked together. Never proceed past a verification gate until it passes.
\end{notebox}

\subsection{Required Tools}

\begin{table}[H]
\centering
\begin{tabularx}{\textwidth}{llX}
\toprule
\textbf{Tool} & \textbf{Purpose} & \textbf{Install} \\
\midrule
Arduino IDE 2.x        & ESP32 firmware          & arduino.cc + Espressif ESP32 board package \\
Python 3.10+           & Pi scripts, ML, backend & Pre-installed on Raspberry Pi OS \\
pymodbus               & MODBUS master on Pi     & \texttt{pip install pymodbus} \\
paho-mqtt              & MQTT client on Pi       & \texttt{pip install paho-mqtt} \\
FastAPI + uvicorn      & REST API + WebSocket    & \texttt{pip install fastapi uvicorn} \\
xgboost, scikit-learn  & ML training/inference   & \texttt{pip install xgboost scikit-learn joblib} \\
Mosquitto              & MQTT broker on Pi       & \texttt{sudo apt install mosquitto} \\
QModMaster             & MODBUS test tool (PC)   & sourceforge.net/projects/qmodmaster \\
VS Code                & IDE for Python + React  & code.visualstudio.com \\
Git                    & Version control         & \texttt{sudo apt install git} \\
\bottomrule
\end{tabularx}
\caption{Development environment tools.}
\end{table}

\subsection{ESP32 Pin Assignment}

\begin{warnbox}
Decide your full pin map \textbf{before} soldering. Changing pins after wires are attached wastes hours.
Get team agreement on this table and treat it as read-only once wiring begins.
\end{warnbox}

\begin{table}[H]
\centering
\small
\begin{tabularx}{\textwidth}{lllX}
\toprule
\textbf{ESP32 Pin} & \textbf{Module} & \textbf{Dir.} & \textbf{Purpose} \\
\midrule
GPIO 34 (ADC1)   & ACS712 OUT      & Input         & Switched current measurement (AC/DC) \\
GPIO 35 (ADC1)   & ZMPT101B OUT    & Input         & Switched voltage measurement (AC) \\
GPIO 21 (SDA)    & INA219 SDA      & Bidirectional & Contact resistance -- I2C data \\
GPIO 22 (SCL)    & INA219 SCL      & Output        & Contact resistance -- I2C clock \\
GPIO 4           & MLX90614 SDA    & Bidirectional & Non-contact temperature sensor \\
GPIO 16 (RX2)    & MAX485 RO       & Input         & MODBUS RTU receive \\
GPIO 17 (TX2)    & MAX485 DI       & Output        & MODBUS RTU transmit \\
GPIO 18          & MAX485 DE/RE    & Output        & RS485 direction control (HIGH = TX) \\
GPIO 2           & Green LED       & Output        & Healthy state \\
GPIO 4           & Yellow LED      & Output        & Warning state \\
GPIO 5           & Red LED         & Output        & Critical / Imminent state \\
GPIO 15          & Buzzer          & Output        & Audio alarm (PWM) \\
GPIO 21/22       & SSD1306 OLED    & Output        & Local display -- HI, cycles, alarm \\
GPIO 36 (ADC)    & Relay coil GPIO & Input         & Relay state detection (switching events) \\
\bottomrule
\end{tabularx}
\caption{ESP32 pin assignment -- finalise before wiring.}
\end{table}

\begin{checkgate}
Upload a blink sketch to the ESP32. Confirm Python imports on the Pi with
\texttt{python -c "import pymodbus, paho.mqtt.client, fastapi"}. Push first Git commit.
\end{checkgate}

% ============================================================
\section{Stage 1 -- IoT Hardware Layer (Day 1)}
% ============================================================

Wire and verify each sensor \textit{individually}. Do not connect multiple sensors until each one
independently produces stable, repeatable readings verified against a reference instrument.

\subsection{Current Sensor -- ACS712}

\textbf{Wiring:} Inline with the relay load circuit. ACS712 needs 5\,V VCC; output must be divided to 3.3\,V
before the ESP32 ADC. Use a 3.3\,k$\Omega$ / 6.8\,k$\Omega$ voltage divider.

\textbf{Verification:} Zero load $\rightarrow$ ADC reads $\approx$2048 (12-bit mid-scale). Known resistive
load $\rightarrow$ reading shifts proportionally. Average 1000 ADC samples per reading to reduce noise.

\begin{warnbox}
ADC input must never exceed 3.3\,V. A direct 5\,V output from ACS712 will damage the ESP32. Always fit
the voltage divider.
\end{warnbox}

\subsection{Voltage Sensor -- ZMPT101B}

\textbf{Wiring:} Connected in parallel with the load (across relay output terminals). Adjust the on-board
potentiometer until output is centred at $\approx$1.65\,V with no load.

\textbf{Verification:} Sample at $\geq$1\,kHz and compute RMS:
$V_{\text{rms}} = \sqrt{\overline{v^2}}$.
Compare to multimeter -- should match within 5\%.

\begin{warnbox}
ZMPT101B outputs a bipolar AC signal centred at 1.65\,V. Negative excursions below 0\,V are clipped by
the ADC -- adjust the offset pot until the waveform is fully within 0--3.3\,V.
\end{warnbox}

\subsection{Temperature Sensor -- MLX90614}

\textbf{Wiring:} I2C at address 0x5A. 4.7\,k$\Omega$ pull-ups on SDA and SCL to 3.3\,V. Aim at relay coil
or contact area (non-contact IR measurement).

\textbf{Verification:} Point at hand vs cold surface -- object temperature changes while ambient stays stable.
Use Adafruit MLX90614 library.

\subsection{Contact Resistance -- INA219 (Most Critical)}

\textbf{Wiring:} 4-wire (Kelvin) measurement. Inject 10\,mA through contacts via a 330\,$\Omega$ resistor
from 3.3\,V. Connect INA219 sense pins directly to contact faces. $R = V_{\text{sense}} / I_{\text{inject}}$.

\textbf{Verification:} New relay: 1--5\,m$\Omega$. Any reading above 50\,m$\Omega$ on a new relay indicates
sense wire resistance -- move clips closer to contact faces.

\begin{warnbox}
The sense wires must carry zero current. Using a single pair of wires for both current injection and voltage
sense adds wire resistance to the measurement. This is the most common contact-resistance mistake.
\end{warnbox}

\subsection{Relay State Detection}

\textbf{Wiring:} 10\,k$\Omega$ pull-down on GPIO. Signal transistor or optocoupler driven by relay coil drive
signal. GPIO HIGH = relay energised.

\textbf{Verification:} Toggle relay manually with a button. GPIO state follows. Implement 5\,ms software
debounce to avoid counting one physical switch as multiple events.

\begin{warnbox}
Without debounce, a single relay closure generates 5--20 false edge detections due to contact bounce,
inflating \texttt{N\_cycles} significantly within hours.
\end{warnbox}

\subsection{Sensor Summary}

\begin{table}[H]
\centering
\begin{tabular}{lllll}
\toprule
\textbf{Sensor} & \textbf{Module} & \textbf{Interface} & \textbf{Range} & \textbf{Verify Against} \\
\midrule
Current       & ACS712-30A     & ADC (GPIO34)      & 0--30 A         & Clamp meter \\
Voltage       & ZMPT101B       & ADC (GPIO35)      & 0--250 V AC     & Multimeter \\
Temperature   & MLX90614       & I2C (0x5A)        & --40 to 125 $^\circ$C & Thermometer \\
Contact-R     & INA219         & I2C (0x40)        & 0--200 m$\Omega$& Milliohm meter \\
Relay State   & GPIO + opto    & Digital (GPIO36)  & HIGH/LOW        & Manual toggle \\
\bottomrule
\end{tabular}
\caption{Sensor summary with verification method.}
\end{table}

\begin{checkgate}
ACS712 reads known current $\pm$5\%. ZMPT101B RMS matches multimeter. MLX90614 responds to hand
proximity. INA219 reads $<$10\,m$\Omega$ for a new relay with correct 4-wire wiring. Relay state GPIO
toggles cleanly with no bounce.
\end{checkgate}

% ============================================================
\section{Stage 2 -- ESP32 Firmware (FreeRTOS) (Day 2)}
% ============================================================

\begin{warnbox}
Do \textbf{not} write a single monolithic \texttt{loop()}. It will block on I2C calls and give inaccurate
timing for cycle counting. Split work across two FreeRTOS tasks pinned to separate cores.
\end{warnbox}

\subsection{Dual-Core Task Split}

\begin{table}[H]
\centering
\begin{tabularx}{\textwidth}{lXX}
\toprule
 & \textbf{Core 0 -- Sensor Task (priority 5)} & \textbf{Core 1 -- Health Task (priority 4)} \\
\midrule
Runs every & 100 ms & Triggered by queue item \\
Does & Read ACS712, ZMPT101B (1000-sample RMS), MLX90614 temp, INA219 resistance, relay GPIO edge with debounce & Compute HI sub-indices and composite HI; run alarm state machine; write MODBUS registers; update OLED (500 ms) \\
Data link & Sends \texttt{SensorData} struct via xQueueSend (non-blocking) & Consumes struct via xQueueReceive \\
NVS writes & Never & Every 100 cycles \\
\bottomrule
\end{tabularx}
\caption{FreeRTOS dual-core task assignment.}
\end{table}

\begin{notebox}
\textbf{Task design rules:}
(1) Core 0 must never be blocked by I2C, MODBUS, or display.
(2) The queue is the only legal shared data path between cores.
(3) Any variable read by both cores needs a \texttt{xSemaphoreCreateMutex()}.
(4) NVS flash has limited write endurance -- write every 100 cycles, not every cycle.
\end{notebox}

\subsection{Firmware Skeleton (Pseudo-code)}

\begin{lstlisting}[language=C, caption={FreeRTOS firmware structure.}]
// Core 0 -- Sensor acquisition, priority 5
void sensor_task(void *param) {
  SensorData d;
  while(1) {
    d.I_rms    = read_acs712_rms(1000);
    d.V_rms    = read_zmpt101b_rms(1000);
    d.T_coil   = mlx.readObjectTempC();
    d.R_cont   = ina219.read_milliohms();
    d.sw_event = detect_relay_edge();   // 5 ms debounce
    xQueueSend(sensor_q, &d, 0);       // non-blocking
    vTaskDelay(pdMS_TO_TICKS(100));
  }
}

// Core 1 -- Health engine, priority 4
void health_task(void *param) {
  SensorData d;
  while(1) {
    xQueueReceive(sensor_q, &d, portMAX_DELAY);
    if (d.sw_event) {
      float E  = d.I_rms * d.I_rms * d.V_rms * T_CLOSE;
      float Kt = arrhenius_factor(d.T_coil);
      W_total += E * Kt;
      N_cycles++;
      if (N_cycles % 100 == 0) nvs_save(N_cycles, W_total);
    }
    float HI = compute_composite_HI(W_total, d.T_coil,
                                    d.R_cont, N_cycles);
    update_alarm_state(HI);
    write_modbus_registers(HI, N_cycles, d);
    update_oled_if_due(HI, N_cycles);   // 500 ms throttle
  }
}
\end{lstlisting}

\begin{checkgate}
OLED shows live HI value. LEDs change colour as you manually change sensor inputs. NVS survives a power
cycle -- N\_cycles and W\_total are restored on reboot.
\end{checkgate}

% ============================================================
\section{Stage 3 -- Health Index Computation (Day 2)}
% ============================================================

The Health Index (HI) is a single number from 0\% (failed) to 100\% (new). It is a weighted composite of
four sub-indices, each capturing a different failure mechanism.

\subsection{Arc Damage Sub-Index}

Arc energy per switching event:
\[
E_{\text{step}} = I^2 \cdot V \cdot t_{\text{close}}
\]

Accumulated with Arrhenius temperature weighting:
\[
W_{\text{total}} = \sum_{\text{events}} \left(E_{\text{step}} \cdot K_T\right)
\qquad
K_T = \exp\!\left[\frac{E_a}{k_B}\!\left(\frac{1}{T_{\text{ref}}} - \frac{1}{T}\right)\right]
\]

\[
HI_{\text{damage}} = 1 - \frac{W_{\text{total}}}{W_{\text{rated}}}
\quad\in[0,1]
\qquad \textbf{Weight: 40\%}
\]

$E_a = 0.7$\,eV, $k_B = 8.617\times10^{-5}$\,eV/K, $T_{\text{ref}} = 298.15$\,K.

\begin{notebox}
A relay switching at 85$^\circ$C accumulates damage 4$\times$ faster than at 25$^\circ$C. This is why
temperature matters even for contact wear -- not just coil insulation.
\end{notebox}

\subsection{Thermal Sub-Index}

\[
HI_{\text{thermal}} = 1 - \frac{T_{\text{coil}}}{T_{\text{max}}}
\quad\in[0,1]
\qquad \textbf{Weight: 25\%}
\]

$T_{\text{max}}$ is from the relay datasheet (typically 85--105$^\circ$C).

\subsection{Contact Resistance Sub-Index}

\[
HI_{\text{resistance}} = 1 - \frac{R_{\text{contact}}}{R_{\text{max}}}
\quad\in[0,1]
\qquad \textbf{Weight: 25\%}
\]

$R_{\text{max}} = 200$\,m$\Omega$ (critical threshold). This sub-index is the most sensitive early-warning
signal -- it responds to contact pitting before HI\_damage or HI\_cycles reach alarming levels.

\subsection{Cycle Sub-Index}

\[
HI_{\text{cycles}} = 1 - \frac{N_{\text{cycles}}}{N_{\text{rated}}}
\quad\in[0,1]
\qquad \textbf{Weight: 10\%}
\]

$N_{\text{rated}}$ from datasheet (e.g.\ 100{,}000 operations).

\subsection{Composite Health Index}

\begin{equation}
\boxed{
HI_{\text{final}} = \bigl(0.40 \cdot HI_{\text{damage}}
                       + 0.25 \cdot HI_{\text{thermal}}
                       + 0.25 \cdot HI_{\text{resistance}}
                       + 0.10 \cdot HI_{\text{cycles}}\bigr) \times 100\%
}
\end{equation}

\subsection{Anomaly Overrides}

Certain sensor events \textbf{bypass HI thresholds} and force an immediate alarm escalation:

\begin{table}[H]
\centering
\begin{tabularx}{\textwidth}{lXl}
\toprule
\textbf{Override trigger} & \textbf{Condition} & \textbf{Forced action} \\
\midrule
Contact resistance spike & $R > 3\times$ baseline & Force alarm = CRITICAL \\
Bounce duration increase& Bounce time $> 1.5\times$ spec & Force alarm = WARNING \\
Thermal runaway         & $dT/dt > 2\,^\circ$C/s  & Force alarm = CRITICAL \\
\bottomrule
\end{tabularx}
\caption{Anomaly overrides -- bypass HI thresholds and escalate alarm directly.}
\end{table}

% ============================================================
\section{Stage 4 -- Alarm State Machine (Day 2)}
% ============================================================

\subsection{Four Alarm States}

\begin{table}[H]
\centering
\renewcommand{\arraystretch}{1.4}
\begin{tabularx}{\textwidth}{llllX}
\toprule
\textbf{State} & \textbf{HI} & \textbf{LED} & \textbf{Buzzer} & \textbf{Action Required} \\
\midrule
Healthy          & $>70\%$   & Solid Green    & Silent        & Log data, normal operation \\
Warning          & 40--70\%  & Solid Yellow   & 1 beep/60\,s  & Alert dashboard + SMS to technician \\
Critical         & 20--40\%  & Slow blink Red & 3 beeps/30\,s & Maintenance work order + email \\
Failure Imminent & $<20\%$   & Fast blink Red & Continuous    & Auto-switch to backup relay \\
\bottomrule
\end{tabularx}
\caption{Alarm state machine -- four states with outputs and required actions.}
\end{table}

\begin{fixbox}
\textbf{Automatic Backup Switching}: when HI drops below 20\%, the ESP32 asserts a GPIO to trigger a
changeover relay \textit{before} failure occurs. Wire a spare GPIO to the backup relay coil through a transistor
driver. This is the difference between predictive maintenance and reactive repair.
\end{fixbox}

\begin{checkgate}
OLED shows live HI. LED changes colour at correct thresholds. Buzzer fires on schedule. NVS restores state
after reboot.
\end{checkgate}

% ============================================================
\section{Stage 5 -- MODBUS Register Map (Day 3)}
% ============================================================

The ESP32 acts as a \textbf{MODBUS RTU slave} at 9600\,baud, 8N1, over RS485 via MAX485.
Use the ArduinoModbus library.

\begin{notebox}
Never change register addresses after the Pi polling script is deployed. Downstream code depends on them.
\end{notebox}

\begin{table}[H]
\centering
\small
\begin{tabularx}{\textwidth}{llllX}
\toprule
\textbf{Addr.} & \textbf{Name} & \textbf{Type} & \textbf{Scale} & \textbf{Description} \\
\midrule
40001 & N\_cycles\_Hi  & UINT16 & $\times$1    & Cycle count -- high 16 bits \\
40002 & N\_cycles\_Lo  & UINT16 & $\times$1    & Cycle count -- low 16 bits \\
40003 & HI\_final      & UINT16 & $\div$1000   & Composite HI (0--1000) \\
40004 & HI\_damage     & UINT16 & $\div$1000   & Damage sub-index \\
40005 & HI\_thermal    & UINT16 & $\div$1000   & Thermal sub-index \\
40006 & HI\_resist     & UINT16 & $\div$1000   & Resistance sub-index \\
40007 & alarm\_state   & UINT16 & $\times$1    & 0=Healthy, 1=Warning, 2=Critical, 3=Imminent \\
40008 & T\_coil        & INT16  & $\div$10     & Temperature in 0.1$^\circ$C units \\
40009 & R\_contact     & UINT16 & $\div$100    & Contact resistance in 0.01\,m$\Omega$ \\
40010 & I\_last\_sw    & UINT16 & $\div$100    & Current at last switch (0.01\,A) \\
40011 & V\_last\_sw    & UINT16 & $\div$10     & Voltage at last switch (0.1\,V) \\
40012 & W\_total       & UINT16 & $\times$1    & Accumulated arc energy (normalised) \\
40013 & dT\_dt         & INT16  & $\div$10     & Temperature rate (0.1$^\circ$C/s) \\
40014 & fault\_bits    & UINT16 & bitmap       & Bit0 = R spike, Bit1 = bounce, Bit2 = thermal runaway \\
\bottomrule
\end{tabularx}
\caption{MODBUS holding register map -- all 4xxxx Holding Registers.}
\end{table}

\subsection{RS485 Bus Rules}

\begin{itemize}[noitemsep]
  \item One RS485 bus supports up to 32 ESP32 slave nodes (one per relay).
  \item Each ESP32 gets a unique MODBUS slave address (1--32) stored in NVS.
  \item Use a 120\,$\Omega$ termination resistor at both ends of the bus for runs $>$3\,m.
  \item Maximum bus length: 1200\,m at 9600\,baud.
  \item The Pi Master polls slaves sequentially: slave 1, slave 2, \ldots{}, slave N, repeat.
\end{itemize}

\begin{checkgate}
Use QModMaster on a laptop (USB-to-RS485 dongle). All 14 registers from every node must be readable and
sane. Values must match the OLED display.
\end{checkgate}

% ============================================================
\section{Stage 6 -- Raspberry Pi Gateway Software (Day 4)}
% ============================================================

The Pi is the bridge between edge nodes and the cloud. Build and verify each module independently before
integrating the next.

\subsection{Pi Software Stack}

\begin{table}[H]
\centering
\begin{tabularx}{\textwidth}{llX}
\toprule
\textbf{Module} & \textbf{Library} & \textbf{Role} \\
\midrule
poll\_relay.py     & pymodbus         & MODBUS master -- polls all 14 registers from each slave \\
sqlite\_store.py   & sqlite3          & Local historian -- appends every poll to time-series DB \\
mqtt\_publish.py   & paho-mqtt + TLS  & Publishes JSON health payloads to broker \\
ml\_inference.py   & xgboost, joblib  & Loads trained model; predicts $\Delta$\,HI per reading \\
fastapi\_server.py & FastAPI, uvicorn & REST API + WebSocket for the dashboard \\
mosquitto          & system service   & MQTT broker on localhost:1883 \\
\bottomrule
\end{tabularx}
\caption{Raspberry Pi software modules.}
\end{table}

\subsection{Step 6.1 -- MODBUS Polling Script}

Build this first. Verify all registers are readable before writing to SQLite.

\begin{lstlisting}[language=Python, caption={MODBUS poll script (core loop).}]
from pymodbus.client import ModbusSerialClient

client = ModbusSerialClient(port="/dev/ttyUSB0", baudrate=9600)
client.connect()

while True:
    for slave_id in range(1, N_SLAVES + 1):
        regs = client.read_holding_registers(40001, 14, slave_id)
        if not regs.isError():
            process_registers(regs.registers, slave_id)
    time.sleep(5)
\end{lstlisting}

\subsection{Step 6.2 -- SQLite Storage}

\begin{itemize}[noitemsep]
  \item Always \texttt{INSERT} (append) -- never overwrite. ML needs full time-series history.
  \item Index on \texttt{(device\_id, ts)} for fast range queries.
  \item Store both raw register values AND computed HI.
  \item Checkpoint WAL every 60 s: \texttt{conn.execute('PRAGMA wal\_checkpoint')}.
\end{itemize}

\subsection{Step 6.3 -- MQTT Publishing}

Publish to topic \texttt{inv/\{device\_id\}/rul} with QoS\,1 and \texttt{retain=True}.

\begin{lstlisting}[language=Python, caption={MQTT JSON payload format (inv/\{id\}/rul).}]
{
  "device_id":   "relay_K1",
  "ts":          1718000000,
  "hi_final":    0.742,
  "hi_damage":   0.810,
  "hi_thermal":  0.920,
  "hi_resistance": 0.680,
  "hi_cycles":   0.851,
  "alarm_state": 0,
  "t_coil":      42.5,
  "r_contact":   8.3,
  "i_last_sw":   12.4,
  "n_cycles":    14823,
  "delta_hi_pred": -0.18,
  "hi_future_pred": 73.9,
  "rul_cycles":  49000
}
\end{lstlisting}

\subsection{Step 6.4 -- FastAPI Backend}

\begin{itemize}[noitemsep]
  \item \texttt{GET /api/fleet} -- current health of all devices (from in-memory dict).
  \item \texttt{GET /api/device/\{id\}/history?limit=200} -- SQLite query.
  \item \texttt{WebSocket /ws} -- push every MQTT message to all connected browsers.
  \item \texttt{POST /api/alerts/\{id\}/acknowledge} -- mark alert handled in DB.
\end{itemize}

\begin{checkgate}
SQLite grows every 5\,s. mosquitto\_sub on Pi shows JSON arriving. /api/fleet returns JSON from browser.
WebSocket push works (test with wscat).
\end{checkgate}

% ============================================================
\section{Stage 7 -- MQTT Topic Architecture (Day 5)}
% ============================================================

\begin{notebox}
Choose topic names carefully -- they are permanent. Changing them later breaks all cloud subscribers.
\end{notebox}

\begin{table}[H]
\centering
\begin{tabularx}{\textwidth}{lllll}
\toprule
\textbf{Topic} & \textbf{Publisher} & \textbf{Payload} & \textbf{QoS} & \textbf{Retain} \\
\midrule
\texttt{inv/\{id\}/rul}     & Pi gateway & Full health JSON  & 1 & Yes \\
\texttt{inv/\{id\}/alarm}   & Pi gateway & Alarm level + fault bitmap & 1 & Yes \\
\texttt{inv/\{id\}/raw}     & Pi gateway & Raw sensor telemetry & 0 & No \\
\texttt{inv/\{id\}/event}   & Pi gateway & Switching events (I, V, ts) & 1 & No \\
\texttt{fleet/summary}      & Cloud      & Aggregated fleet health & 1 & Yes \\
\texttt{edge/\{id\}/config} & Cloud      & OTA threshold update & 1 & No \\
\bottomrule
\end{tabularx}
\caption{MQTT topic structure.}
\end{table}

% ============================================================
\section{Stage 8 -- ML Pipeline: Delta-HI XGBoost (Days 6--7)}
% ============================================================

\subsection{The Problem with the Old Model (Absolute HI Target)}

\begin{warnbox}
\textbf{v1/v2 failure mode (do not repeat):} The original model predicted absolute future HI
(\texttt{target = HI\_future}). Training data covered HI 34--99\% (healthy relay). Test data covered
HI 5--44\% (degraded relay). The model had never seen low-HI outputs during training, so it defaulted to
predicting the training mean ($\approx$40\%) for every input -- a flat, useless line.

This is the classic \textbf{covariate shift} problem.
\end{warnbox}

\subsection{The Fix: Predict Delta-HI}

\begin{fixbox}
\textbf{Core change in v3:} predict the \textit{change} in HI, not the absolute future value.
\[
\Delta HI = HI[t + \texttt{FUTURE\_STEPS}] - HI[t]
\]
$\Delta HI$ is always small ($-3$ to $+3$ percentage points) regardless of whether the relay is healthy or
degraded. There is \textbf{no covariate shift} -- the output distribution is the same throughout the lifecycle.

Reconstruct absolute future HI trivially at inference:
\[
\widehat{HI}_{\text{future}} = HI_{\text{current}} + \widehat{\Delta HI}
\]
\end{fixbox}

\begin{table}[H]
\centering
\begin{tabularx}{\textwidth}{lXX}
\toprule
\textbf{Property} & \textbf{Old: absolute HI target} & \textbf{New: $\Delta$HI target (v3)} \\
\midrule
Output range    & 0--100\% (varies with relay state)  & $-3$ to $+3$\% (nearly constant) \\
Covariate shift & Severe                              & None \\
Generalisation  & Poor on unseen HI ranges            & Strong across full lifecycle \\
Reconstruction  & Direct output                       & $HI_{\text{now}} + \Delta HI$ \\
\bottomrule
\end{tabularx}
\caption{Old vs.\ new prediction target comparison.}
\end{table}

\subsection{Training Data: Physics-Based Multi-Stage Degradation}

Run \texttt{generate\_sample\_data.py} to produce \texttt{relay\_data.xlsx} (3000 rows):

\begin{table}[H]
\centering
\begin{tabular}{llll}
\toprule
\textbf{Phase} & \textbf{Rows} & \textbf{Degradation} & \textbf{What it models} \\
\midrule
Burn-in   & 0--300   & 0.00--0.12 & Quick initial settling \\
Normal    & 300--2250 & 0.12--0.52 & Slow, steady linear wear \\
Wear-out  & 2250--3000 & 0.52--1.00 & Accelerating contact erosion \\
\bottomrule
\end{tabular}
\caption{Three-phase bathtub degradation curve for training data.}
\end{table}

The generator also injects \textbf{25 overload events} (current $\times$1.8--3.5) and
\textbf{12 hot spells} ($+20^\circ$C for 30 readings) to ensure the model learns realistic stress events.
Because all phases and events appear throughout the dataset, an 80/20 chronological split always includes
both healthy and degraded examples in both halves.

\subsection{Feature Engineering}

\begin{table}[H]
\centering
\begin{tabularx}{\textwidth}{lX}
\toprule
\textbf{Feature Group} & \textbf{Features} \\
\midrule
Base sensors     & current, voltage, temp, cycles, resistance, peak\_current, frequency \\
Derived sensor   & dT\_dt (finite difference of temperature) \\
Physics outputs  & HI, HI\_damage, HI\_thermal, HI\_resistance, HI\_cycles \\
Cumulative damage & damage ($W_{\text{total}}$, thermally weighted cumulative arc energy) \\
Lifecycle        & lifecycle\_pos $= t / t_{\max}$ (where in life the relay is) \\
HI lags          & HI\_lag\_k for $k \in \{1,2,3,5,10,20\}$ \\
HI slopes        & Least-squares slope of HI over windows $w \in \{5,10,20,50\}$ \\
HI rolling stats & Rolling mean \& std of HI over $w \in \{5,10,20,50\}$ \\
Sensor rolling   & Rolling mean \& std of current, temp, resistance, voltage over $w \in \{5,10,20,50\}$ \\
Interactions     & damage $\times$ temp; current $\times$ resistance; HI $\times$ cycles; HI slope $\times$ damage \\
\bottomrule
\end{tabularx}
\caption{Feature engineering -- 60--80 columns total.}
\end{table}

\subsection{XGBoost Hyperparameters}

\begin{table}[H]
\centering
\begin{tabular}{lll}
\toprule
\textbf{Parameter} & \textbf{Value} & \textbf{Rationale} \\
\midrule
n\_estimators           & 2000   & Large pool; early stopping finds optimal \\
max\_depth              & 4      & Shallow trees -- reduces overfitting on noisy $\Delta$HI \\
learning\_rate          & 0.02   & Slow, stable convergence \\
subsample               & 0.8    & Row sampling -- adds stochasticity \\
colsample\_bytree       & 0.8    & Feature sampling per tree \\
min\_child\_weight      & 10     & Smoother predictions \\
gamma                   & 0.2    & Minimum split gain \\
reg\_alpha / reg\_lambda & 0.1 / 2.0 & L1 + L2 regularisation \\
early\_stopping\_rounds & 50     & Stops when val-RMSE does not improve \\
\bottomrule
\end{tabular}
\caption{XGBoost hyperparameters for $\Delta$HI regression.}
\end{table}

\begin{notebox}
\texttt{RobustScaler} (scikit-learn) is used instead of \texttt{MinMaxScaler} because it is based on
median and IQR rather than min/max, making it resistant to the overload current spikes present in the data.
\end{notebox}

\subsection{Train/Test Split}

80\% / 20\% chronological split (no shuffle). The first 80\% of time-ordered rows train the model; the last
20\% are the held-out test set. Because we predict $\Delta$HI rather than absolute HI, the test distribution
matches training regardless of HI level.

\subsection{Training Pipeline (Google Colab)}

\textbf{Preferred training environment}: Google Colab (\href{https://colab.research.google.com}{colab.research.google.com}). The script \texttt{RelayGuard\_Colab.py} is already organised as Colab cells.

\begin{enumerate}[noitemsep]
  \item Open Colab. Create a new notebook.
  \item Paste each cell from \texttt{RelayGuard\_Colab.py} into a separate code cell.
  \item Run Cell 3: upload \texttt{relay\_data.xlsx} (generated by \texttt{generate\_sample\_data.py}) or set \texttt{USE\_SAMPLE\_DATA = True}.
  \item Run cells top to bottom. Training completes in $\approx$2--5 minutes.
  \item Cell 8 downloads: \texttt{relay\_guard\_model.pkl}, \texttt{scaler.pkl},
        \texttt{predictions.xlsx}, \texttt{training\_results.png}.
\end{enumerate}

\subsection{Core Training Code}

\begin{lstlisting}[language=Python, caption={Delta-HI target creation and XGBoost training (from RelayGuard\_Colab.py).}]
# -- Delta HI target (the key fix) ------------------------
df["HI_future"] = df["HI"].shift(-FUTURE_STEPS)
df["delta_HI"]  = df["HI_future"] - df["HI"]
df.dropna(subset=["delta_HI"], inplace=True)

# -- Features / target split -----------------------------
EXCLUDE   = {"HI_future", "delta_HI", "Alert", "time"}
feat_cols = [c for c in df.columns if c not in EXCLUDE]
X, y      = df[feat_cols].values, df["delta_HI"].values

# -- Outlier-resistant scaling ---------------------------
scaler   = RobustScaler()
X_scaled = scaler.fit_transform(X)

# -- Time-ordered 80/20 split ----------------------------
split       = int(len(X) * 0.80)
X_tr, X_te  = X_scaled[:split], X_scaled[split:]
y_tr, y_te  = y[:split], y[split:]
hi_te       = df["HI"].values[split:]   # current HI

# -- Train -----------------------------------------------
model = XGBRegressor(**XGB_PARAMS)
model.fit(X_tr, y_tr,
          eval_set=[(X_tr, y_tr), (X_te, y_te)],
          verbose=100)

# -- Reconstruct absolute HI from delta ------------------
delta_pred = model.predict(X_te)
hi_pred    = (hi_te + delta_pred).clip(0, 100)
\end{lstlisting}

\subsection{Evaluation Targets}

\begin{table}[H]
\centering
\begin{tabular}{lll}
\toprule
\textbf{Metric} & \textbf{Formula} & \textbf{Target} \\
\midrule
MAE   & $\frac{1}{n}\sum|HI_{\text{actual}} - HI_{\text{pred}}|$ & $< 3\%$ \\
RMSE  & $\sqrt{\frac{1}{n}\sum(HI_{\text{actual}} - HI_{\text{pred}})^2}$ & $< 5\%$ \\
R$^2$ & $1 - \text{SS}_\text{res}/\text{SS}_\text{tot}$ & $> 0.90$ \\
\bottomrule
\end{tabular}
\caption{ML evaluation targets (evaluated on reconstructed absolute HI, not $\Delta$HI).}
\end{table}

\subsection{LSTM Anomaly Detection (Advanced)}

Beyond XGBoost, the full system design includes an LSTM autoencoder for anomaly detection:

\begin{itemize}[noitemsep]
  \item \textbf{Input}: sliding window of 50 timesteps $\times$ 10 features (HI sub-indices, T, R, I, dT/dt, resistance slope, switching frequency).
  \item \textbf{Architecture}: Encoder LSTM(32) $\to$ bottleneck $\to$
        Decoder LSTM(32) $\to$ TimeDistributed Dense(10).
  \item \textbf{Target}: reconstruct the input sequence (autoencoder -- no failure labels needed).
  \item \textbf{Anomaly}: reconstruction error $>$ mean $+ 3\sigma$ on healthy training set.
  \item \textbf{Value}: fires 10--20 data points \textit{before} any rule-based threshold triggers. This is the AI early-warning advantage.
  \item Train on Google Colab GPU. Export to TFLite. Deploy \texttt{.tflite} to Pi for inference.
\end{itemize}

\subsection{Weibull Survival Analysis (Probabilistic RUL)}

Weibull gives a probability distribution over remaining life rather than a point estimate:

\begin{align}
\text{RUL}_{50\%} &= \eta \cdot (\ln 2)^{1/\beta} \\
\text{RUL}_{85\%} &= \eta \cdot (-\ln 0.15)^{1/\beta}
\end{align}

Parameters: $\beta \approx 2.5$--3.5 (wear-out failure, increasing hazard), $\eta$ fitted from B10 life data
or actual field failures.

\begin{checkgate}
\texttt{ml\_inference.py} loads \texttt{relay\_guard\_model.pkl} and runs inference without errors.
Predicted $\Delta$HI values are in the range $-3$ to $+3$. Anomaly flag appears in API output.
\end{checkgate}

\subsection{Deploying the Model on the Pi}

\begin{lstlisting}[language=Python, caption={Inference on Raspberry Pi using saved model.}]
import joblib, numpy as np

model  = joblib.load("relay_guard_model.pkl")
scaler = joblib.load("scaler.pkl")

# Build feature vector using same physics + feature
# engineering pipeline as training, then:
X_new      = scaler.transform(df_new[feat_cols].values)
delta_hi   = model.predict(X_new)[0]

hi_current = df_new["HI"].iloc[-1]
hi_future  = float(np.clip(hi_current + delta_hi, 0, 100))

print(f"Current HI : {hi_current:.1f}%")
print(f"Predicted  : {hi_future:.1f}%")
print(f"Alert      : {classify_alert(hi_future)}")
\end{lstlisting}

% ============================================================
\section{Stage 9 -- Dashboard Development (Day 6)}
% ============================================================

Build the dashboard \textbf{after} the data pipeline is live. Real sensor data makes the demo compelling;
mocked data makes it unconvincing.

\subsection{Dashboard Panels (Priority Order)}

\begin{table}[H]
\centering
\begin{tabularx}{\textwidth}{llX}
\toprule
\textbf{Priority} & \textbf{Panel} & \textbf{Contents} \\
\midrule
P1 & Fleet Overview Grid   & One card per relay: ID, large HI\% coloured green/amber/red, alarm state, RUL estimate. \\
P2 & HI Timeline           & Line chart (last 7 days) of HI\_final + sub-index lines. Vertical markers at alarm changes. Red dots = LSTM anomaly fires. \\
P3 & Contact Resistance Trend & R\_contact in m$\Omega$ over time. Horizontal lines at 50\,m$\Omega$ (warning) and 200\,m$\Omega$ (critical). \\
P4 & Current vs Damage Scatter & Switching events: I\_switch on x-axis, arc energy on y-axis, coloured by temperature. \\
P5 & Work Order Panel      & Auto-populates when alarm\_state $\geq$\,2. Shows relay ID, HI, predicted failure date, acknowledge button. \\
Bonus & Weibull Survival Curve & Probability of survival vs remaining cycles. 50th, 85th, 95th percentile bands. \\
\bottomrule
\end{tabularx}
\caption{Dashboard panels in build priority order.}
\end{table}

\subsection{Technology Options}

\begin{table}[H]
\centering
\begin{tabularx}{\textwidth}{lXll}
\toprule
\textbf{Option} & \textbf{Stack} & \textbf{Setup time} & \textbf{Best for} \\
\midrule
React + Recharts   & React, Recharts, WebSocket client & 2--3 days   & Custom UI, demo polish \\
Grafana + InfluxDB & Grafana OSS, InfluxDB v2, Telegraf & 2--4 hours & Quick dashboard, no-code \\
Node-RED Dashboard & Node-RED + node-red-dashboard     & 1--2 hours  & Fastest prototype \\
\bottomrule
\end{tabularx}
\caption{Dashboard technology options.}
\end{table}

\begin{checkgate}
HI gauge updates in browser without page refresh. Alarm state change visible within 10\,s. All 5 panels
populated with live data.
\end{checkgate}

% ============================================================
\section{Stage 10 -- Integration and Demo (Day 7)}
% ============================================================

\subsection{Full Integration Test Sequence}

Run the full end-to-end pipeline manually:

\begin{enumerate}[noitemsep]
  \item Power on ESP32 nodes. Confirm OLED shows live HI.
  \item Run Pi poll script. Confirm SQLite grows every 5\,s.
  \item Confirm MQTT JSON arrives on mosquitto\_sub.
  \item Open dashboard in browser. Confirm all panels update live.
  \item Simulate stress: apply overload current to relay. Confirm HI drops and alarm escalates on dashboard.
  \item Confirm work order card auto-populates at alarm\_state 2.
  \item Confirm backup relay GPIO fires when HI $<$ 20\%.
\end{enumerate}

\subsection{Demo Script (3 Minutes)}

\begin{table}[H]
\centering
\begin{tabularx}{\textwidth}{llX}
\toprule
\textbf{Step} & \textbf{What you do} & \textbf{What the audience sees} \\
\midrule
1 & Start with all relays healthy & Dashboard all green. ESP32 LED = steady green. \\
2 & Increase stress (potentiometer or injected overload) & HI drops live on dashboard. \\
3 & Reach Warning ($\approx$60\%) & Yellow LED. LSTM anomaly dot appears on timeline \textit{before} rule threshold fires. Emphasise this gap. \\
4 & Reach Critical ($\approx$30\%) & Red LED + slow beep. Work order card pops up: ``Replace Relay K1 -- failure in 2.1 hours''. \\
5 & Reach Failure Imminent ($<$20\%) & Fast blink + continuous buzzer. Dashboard red. Backup relay switches automatically. \\
6 & Business case closing & ``Unplanned downtime costs X/hour. This system predicted failure 48\,h early, enabling a planned replacement.'' \\
\bottomrule
\end{tabularx}
\caption{Demo scenario script.}
\end{table}

\begin{notebox}
\textbf{Always record a backup video on Day 5.} If WiFi fails or hardware misbehaves on demo day, play the
video without breaking stride. Judges care about the concept and build quality -- a backup demonstrates
professionalism.
\end{notebox}

% ============================================================
\section{Known Issues and Design Decisions}
% ============================================================

\begin{table}[H]
\centering
\begin{tabularx}{\textwidth}{lXl}
\toprule
\textbf{Issue} & \textbf{Resolution} & \textbf{Status} \\
\midrule
Flat $\approx$40\% prediction     & Switched to $\Delta$HI target (eliminates covariate shift)         & Fixed in v3 \\
MinMaxScaler distorted by spikes  & Replaced with RobustScaler (IQR-based)                            & Fixed in v3 \\
Training data: healthy phase only & Multi-stage bathtub curve covers all phases in every split         & Fixed (data gen v2) \\
train\_model.py uses absolute HI  & Needs updating to match Colab v3 $\Delta$HI logic                 & Planned \\
RUL estimate is linear            & Adequate short-term; learned RUL head would improve long-term RUL & Future work \\
No real relay data yet            & Physics + bathtub simulation used as surrogate                     & Ongoing \\
LCD / MOSFET / switch models      & Defined in RUL doc but not yet implemented in firmware             & Future work \\
\bottomrule
\end{tabularx}
\caption{Known issues, resolutions, and status.}
\end{table}

% ============================================================
\section{Production Roadmap}
% ============================================================

\subsection{Phase 1 -- Edge Instrumentation (Months 1--3)}
\begin{itemize}[noitemsep]
  \item Instrument firmware with cycle counters and stress-weighted wear models for all 5 component types.
  \item Validate degradation models against accelerated life test data.
  \item Implement non-volatile health register storage.
  \item Expose all data over MODBUS at agreed register map.
\end{itemize}

\subsection{Phase 2 -- Connectivity (Months 3--5)}
\begin{itemize}[noitemsep]
  \item Deploy MODBUS polling gateway + MQTT pipeline.
  \item Stand up InfluxDB / TimescaleDB time-series database.
  \item Build minimal fleet dashboard (RUL gauge per unit).
\end{itemize}

\subsection{Phase 3 -- Cloud ML (Months 5--9)}
\begin{itemize}[noitemsep]
  \item Collect 6+ months of fleet telemetry.
  \item Train LSTM autoencoder and refine XGBoost $\Delta$HI model on real data.
  \item Deploy Weibull survival models fitted to actual field failures.
  \item Enable OTA threshold updates from cloud back to edge devices.
\end{itemize}

\subsection{Phase 4 -- Integration and Scale (Months 9--11)}
\begin{itemize}[noitemsep]
  \item Connect work-order triggers to CMMS (SAP PM, IBM Maximo, or webhooks).
  \item Build fleet-wide RUL heatmap.
  \item Add confidence intervals to RUL estimates.
  \item Establish automated monthly retraining pipeline with new field data.
\end{itemize}

% ============================================================
\section{Team Responsibility Matrix}
% ============================================================

\begin{table}[H]
\centering
\begin{tabularx}{\textwidth}{lllX}
\toprule
\textbf{Role} & \textbf{Stages} & \textbf{Day} & \textbf{Deliverable} \\
\midrule
Hardware Lead   & 0, 1, 2, 3 & Day 3 & All sensors verified, MODBUS readable from laptop \\
Backend Lead    & 4, 5, 6    & Day 5 & Pi pipeline live, API serving JSON, WebSocket push working \\
ML Lead         & 8          & Day 6 & XGBoost deployed on Pi, inference returns delta HI \\
Frontend Lead   & 9          & Day 6 & Dashboard live with real data, all 5 panels visible \\
Full Team       & 10         & Day 7 & End-to-end test, demo rehearsal, backup video recorded \\
\bottomrule
\end{tabularx}
\caption{Team responsibility matrix.}
\end{table}

% ============================================================
\section{Summary -- Complete Project Checklist}
% ============================================================

\begin{tcolorbox}[colback=notebg,colframe=accent,title={\bfseries RelayGuard -- Complete Build Checklist}]
\begin{enumerate}[noitemsep]
  \item \textbf{Day 0:} Tools installed. Pin map agreed. Git repo created.
  \item \textbf{Day 1:} Each sensor (ACS712, ZMPT101B, MLX90614, INA219, relay GPIO) independently verified.
  \item \textbf{Day 2:} FreeRTOS firmware running. OLED shows live HI. LEDs respond. NVS survives reboot.
  \item \textbf{Day 3:} MODBUS slave on ESP32. QModMaster reads all 14 registers from Pi.
  \item \textbf{Day 4:} Pi gateway polling. SQLite growing. MQTT JSON on broker.
  \item \textbf{Day 5:} FastAPI backend live. WebSocket push working. Full backup video recorded.
  \item \textbf{Day 6:} React/Grafana dashboard live with all 5 panels. XGBoost $\Delta$HI model deployed on Pi.
  \item \textbf{Day 7:} End-to-end system test. Stress scenario drives full alarm journey (green to red). Work order auto-creates. Demo rehearsed.
  \item \textbf{ML:} Target MAE $< 3\%$, RMSE $< 5\%$, R$^2 > 0.90$ on reconstructed absolute HI.
  \item \textbf{Architecture:} Predict $\Delta$HI (not absolute HI) to eliminate covariate shift.
\end{enumerate}
\end{tcolorbox}

\vspace{1.5em}
\hrule
\vspace{0.5em}
{\small\centering
\textit{RelayGuard Complete Project Development Guide}\quad|\quad Version 3.0\quad|\quad April 2026\\
Supersedes: Relay\_Health\_System\_Dev\_Guide.pdf, RUL\_Predictive\_Maintenance\_Solution.pdf, ML\_Model.pdf
\par}

\end{document}
